% GENERATED FILE. Edit canonical lessons, then run npm run generate:books.
% canonical-source-hash: fnv1a64:9b8e922c93a5159f
% canonical-lessons: PT-C494-fralda, PT-C494-biberao, PT-C494-chupeta, PT-C494-gatinhar, PT-C494-embalar

\chapter{Fraldas e biberões}
\label{ch:pt-fraldas-e-biberoes}
\hlchaptermodality{494}

\begin{chapteropening}
\textbf{By the end of this chapter:} \emph{I can talk about looking after a baby.}

Complete the last lesson of chapter 494: I can talk about looking after a baby.
\end{chapteropening}

\section[a fralda]{a fralda — nappy, diaper}

\label{lesson:PT-C494-fralda}



\begin{quote}
Before the new one: say the Portuguese for to adopt, then the Portuguese for adoption.
\end{quote}

\subsection*{You'll want to know: a fralda}
\textbf{a fralda} — \textquotedblleft{}nappy, diaper\textquotedblright{}.

\begin{quote}
O pai mudou a fralda ao bebé.
\end{quote}

\begin{quote}
\emph{The father changed the baby's nappy.}
\end{quote}

\emph{Mudar a fralda} — to change a nappy.

\begin{cousinweb}
\emph{fralda} goes back to Old Portuguese \emph{faldra} (EtymDB), from a Gothic word related to \emph{falþan} \textquotedblleft{}to fold\textquotedblright{}.
\end{cousinweb}

\begin{grammarlens}[title={What you've built}]
One more word for telling, explaining and describing at B1.
\end{grammarlens}

\subsection*{Your turn}
\begin{itemize}
\raggedright
  \item \emph{Say it:} \emph{a fralda}
  \item \emph{Say it:} \emph{a fralda}, once more
  \item \emph{Say it:} a sentence of your own with \emph{a fralda}
  \item \emph{Recall:} say the Portuguese for to adopt, then the Portuguese for adoption, then say \emph{a fralda} again
\end{itemize}

\begin{tcolorbox}[breakable,colback=teal!4,colframe=teal!35!black,title={Before you move on}]
What does a fralda mean? (\textquotedblleft{}Nappy, diaper\textquotedblright{}.) Say it once more.
\end{tcolorbox}
\section[o biberão]{o biberão — baby's bottle}

\label{lesson:PT-C494-biberao}



\begin{quote}
Before the new one: say the Portuguese for adoption, then the Portuguese for nappy.
\end{quote}

\subsection*{You'll want to know: o biberão}
\textbf{o biberão} — \textquotedblleft{}baby's bottle\textquotedblright{}.

\begin{quote}
A avó aqueceu o leite do biberão.
\end{quote}

\begin{quote}
\emph{Grandma warmed the milk in the baby's bottle.}
\end{quote}

In Brazil: \emph{a mamadeira}.

\begin{cousinweb}
\emph{biberão} was borrowed from French \emph{biberon} (EtymDB), built on Latin \emph{bibere} \textquotedblleft{}to drink\textquotedblright{}.
\end{cousinweb}

\begin{grammarlens}[title={What you've built}]
One more word for telling, explaining and describing at B1.
\end{grammarlens}

\subsection*{Your turn}
\begin{itemize}
\raggedright
  \item \emph{Say it:} \emph{o biberão}
  \item \emph{Say it:} \emph{o biberão}, once more
  \item \emph{Say it:} a sentence of your own with \emph{o biberão}
  \item \emph{Recall:} say the Portuguese for adoption, then the Portuguese for nappy, then say \emph{o biberão} again
\end{itemize}

\begin{tcolorbox}[breakable,colback=teal!4,colframe=teal!35!black,title={Before you move on}]
What does o biberão mean? (\textquotedblleft{}Baby's bottle\textquotedblright{}.) Say it once more.
\end{tcolorbox}
\section[a chupeta]{a chupeta — dummy, pacifier}

\label{lesson:PT-C494-chupeta}



\begin{quote}
Before the new one: say the Portuguese for nappy, then the Portuguese for baby's bottle.
\end{quote}

\subsection*{You'll want to know: a chupeta}
\textbf{a chupeta} — \textquotedblleft{}dummy, pacifier\textquotedblright{}.

\begin{quote}
O bebé só adormece com a chupeta.
\end{quote}

\begin{quote}
\emph{The baby only falls asleep with the dummy.}
\end{quote}

Built on \emph{chupar} \textquotedblleft{}to suck\textquotedblright{}.

\begin{cousinweb}
\emph{chupeta} is built on \emph{chupar} \textquotedblleft{}to suck\textquotedblright{} (EtymDB); the verb is usually taken to imitate the sound of sucking.
\end{cousinweb}

\begin{grammarlens}[title={What you've built}]
One more word for telling, explaining and describing at B1.
\end{grammarlens}

\subsection*{Your turn}
\begin{itemize}
\raggedright
  \item \emph{Say it:} \emph{a chupeta}
  \item \emph{Say it:} \emph{a chupeta}, once more
  \item \emph{Say it:} a sentence of your own with \emph{a chupeta}
  \item \emph{Recall:} say the Portuguese for nappy, then the Portuguese for baby's bottle, then say \emph{a chupeta} again
\end{itemize}

\begin{tcolorbox}[breakable,colback=teal!4,colframe=teal!35!black,title={Before you move on}]
What does a chupeta mean? (\textquotedblleft{}Dummy, pacifier\textquotedblright{}.) Say it once more.
\end{tcolorbox}
\section[gatinhar]{gatinhar — to crawl (like a baby)}

\label{lesson:PT-C494-gatinhar}



\begin{quote}
Before the new one: say the Portuguese for baby's bottle, then the Portuguese for dummy.
\end{quote}

\subsection*{You'll want to know: gatinhar}
\textbf{gatinhar} — \textquotedblleft{}to crawl (like a baby)\textquotedblright{}.

\begin{quote}
Com oito meses, a bebé já gatinha pela casa.
\end{quote}

\begin{quote}
\emph{At eight months, the baby is already crawling round the house.}
\end{quote}

Built on \emph{gato}: to move on all fours like a little cat.

\begin{cousinweb}
\emph{gatinhar} is built on \emph{gatinho} \textquotedblleft{}kitten\textquotedblright{}, from \emph{gato} \textquotedblleft{}cat\textquotedblright{}, which is inherited from Late Latin \emph{cattus} (EtymDB).
\end{cousinweb}

\begin{grammarlens}[title={What you've built}]
One more word for telling, explaining and describing at B1.
\end{grammarlens}

\subsection*{Your turn}
\begin{itemize}
\raggedright
  \item \emph{Say it:} \emph{gatinhar}
  \item \emph{Say it:} \emph{gatinhar}, once more
  \item \emph{Say it:} a sentence of your own with \emph{gatinhar}
  \item \emph{Recall:} say the Portuguese for baby's bottle, then the Portuguese for dummy, then say \emph{gatinhar} again
\end{itemize}

\begin{tcolorbox}[breakable,colback=teal!4,colframe=teal!35!black,title={Before you move on}]
What does gatinhar mean? (\textquotedblleft{}To crawl (like a baby)\textquotedblright{}.) Say it once more.
\end{tcolorbox}
\section[embalar]{embalar — to rock (a baby); to pack}

\label{lesson:PT-C494-embalar}



\begin{quote}
Before the new one: say the Portuguese for dummy, then the Portuguese for to crawl.
\end{quote}

\subsection*{You'll want to know: embalar}
\textbf{embalar} — \textquotedblleft{}to rock (a baby); to pack\textquotedblright{}.

\begin{quote}
A mãe embalou o bebé até ele adormecer.
\end{quote}

\begin{quote}
\emph{The mother rocked the baby until he fell asleep.}
\end{quote}

It also means to pack goods: \emph{embalar as compras}.

\begin{cousinweb}
In the sense \textquotedblleft{}to pack\textquotedblright{}, \emph{embalar} comes from French \emph{emballer} (EtymDB), built on \emph{balle} \textquotedblleft{}a bale\textquotedblright{}; the origin of the sense \textquotedblleft{}to rock a baby\textquotedblright{} is not settled.
\end{cousinweb}

\begin{grammarlens}[title={What you've built}]
One more word for telling, explaining and describing at B1.
\end{grammarlens}

\subsection*{Your turn}
\begin{itemize}
\raggedright
  \item \emph{Say it:} \emph{embalar}
  \item \emph{Say it:} \emph{embalar}, once more
  \item \emph{Say it:} a sentence of your own with \emph{embalar}
  \item \emph{Recall:} say the Portuguese for dummy, then the Portuguese for to crawl, then say \emph{embalar} again
\end{itemize}

\begin{tcolorbox}[breakable,colback=teal!4,colframe=teal!35!black,title={Before you move on}]
What does embalar mean? (\textquotedblleft{}To rock (a baby); to pack\textquotedblright{}.) Say it once more.
\end{tcolorbox}
